\documentclass[a4paper,11pt]{article}

% Paquetes básicos
\usepackage[utf8]{inputenc}
\usepackage[T1]{fontenc}
\usepackage[spanish]{babel}
\usepackage{amsmath, amssymb}
\usepackage{bm}
\usepackage{geometry}
\geometry{margin=2.5cm}

\title{Modelo binivel con PPA porcentaje constante anual}
\author{}
\date{}

\begin{document}

\maketitle

\section{Nivel inferior: clearing del mercado eléctrico}
\subsection{Definición del nivel inferior}

Sean:
\begin{itemize}
  \item $J$: conjunto de generadores, con $i \in J$ el agente estratégico renovable.
  \item $L$: conjunto de consumidores.
  \item $T$: conjunto de horas de mercado.
  \item $Y$: conjunto de años naturales, con mapeo $y:T\to Y$.
\end{itemize}

Parámetros:
\begin{itemize}
  \item $\alpha_{j,t}$: coste marginal del generador $j$ en hora $t$.
  \item $\beta_{l,t}$: utilidad marginal de demanda $l$ en hora $t$.
  \item $Q_{j,t}$: capacidad máxima ofertable del generador $j$.
  \item $D_{l,t}$: demanda máxima de $l$ en hora $t$.
  \item $x_{i,y} \in [0,1]$: fracción de producción del agente $i$ vendida como PPA en el año $y$ (decisión del nivel superior).
  \item $\gamma$: precio del PPA
\end{itemize}

Problema del nivel inferior:

\begin{align}
\max_{q,d}\;\;
& \sum_{l,t} \beta_{l,t}\, d_{l,t}
\;-\; \sum_{j,t} \alpha_{j,t}\, q_{j,t}
\label{LF_obj}\\[2mm]
\text{s.a.}\;\;
& \sum_{j} q_{j,t} \;-\; \sum_{l} d_{l,t}  \;=\; 0 
\qquad (\perp\; \lambda_t)\quad \forall t
\label{LF_balance}\\
& 0 \le q_{j,t} \le Q_{j,t} 
\quad (\perp\; \omega^{q,\min}_{j,t},\;\omega^{q,\max}_{j,t})
\quad \forall j\neq i,\; \forall t
\label{LF_cap_others}\\
& 0 \le q_{i,t} \le (1-x_{i,y(t)})\,Q_{i,t}
\quad (\perp\; \omega^{q,\min}_{i,t},\;\omega^{q,\max}_{i,t})
\quad \forall t
\label{LF_cap_i}\\
& 0 \le d_{l,t} \le D_{l,t} 
\quad (\perp\; \omega^{d,\min}_{l,t},\;\omega^{d,\max}_{l,t})
\quad \forall l,\; \forall t.
\label{LF_dem}
\end{align}

\subsection{Lagrangiano del problema inferior}

El Lagrangiano asociado al problema de clearing del mercado eléctrico es:

\begin{align}
\mathcal{L} =\;
& \sum_{l,t} \beta_{l,t}\, d_{l,t}
- \sum_{j,t} \alpha_{j,t}\, q_{j,t} \\
& + \sum_{t} \lambda_t \left( \sum_{j} q_{j,t} - \sum_{l} d_{l,t} \right) \\
& + \sum_{j \neq i,\, t} \omega^{q,\min}_{j,t}\, q_{j,t}
+ \sum_{j \neq i,\, t} \omega^{q,\max}_{j,t}\, (Q_{j,t} - q_{j,t}) \\
& + \sum_{t} \omega^{q,\min}_{i,t}\, q_{i,t}
+ \sum_{t} \omega^{q,\max}_{i,t}\, \left( (1-x_{i,y(t)}) Q_{i,t} - q_{i,t} \right) \\
& + \sum_{l,t} \omega^{d,\min}_{l,t}\, d_{l,t}
+ \sum_{l,t} \omega^{d,\max}_{l,t}\, (D_{l,t} - d_{l,t})
\end{align}

Importante apuntar que es un problema de maximizacion y esto cambia el signo de las restricciones de desigualdad al introducirlas en el Lagrangiano.

\subsection{KKT del nivel inferior}
\begin{itemize}
    \item \textbf{Factibilidad primal:}
    Son las restricciones originales del problema.
    \begin{align}
        &\forall t: && \sum_{j} q_{j,t} - \sum_{l} d_{l,t} = 0 \\
        &\forall j \neq i,\, \forall t: && 0 \leq q_{j,t} \leq Q_{j,t} \\
        &\forall t: && 0 \leq q_{i,t} \leq (1-x_{i,y(t)}) Q_{i,t} \\
        &\forall l,\, \forall t: && 0 \leq d_{l,t} \leq D_{l,t}
    \end{align}

    \item \textbf{Condiciones de complementariedad:}
    Para cada restricción la variable no se encuentra en el límite (la variable del dual es 0) o bien al encontrarse en el límite dicha restricción puede ser responsable de no poder mejorar el valor de la función objetivo (la variable de dual puede no ser 0). Importante notar que no se aplica condición de complementariedad para las restricciones de igualdad, aunque estuviera indicado en la descripción del problema, pues esto es para indicar cual es la variable dual asociada para poder nombrarla en el nivel superior.
    \begin{align}
        &\forall j,\, t: && \omega^{q,\min}_{j,t} \cdot q_{j,t} = 0 \\
        &\forall j \neq i,\, t: && \omega^{q,\max}_{j,t} \cdot (q_{j,t} - Q_{j,t}) = 0 \\
        &\forall t: && \omega^{q,\max}_{i,t} \cdot (q_{i,t} - (1-x_{i,y(t)}) Q_{i,t}) = 0 \\
        &\forall l,\, t: && \omega^{d,\min}_{l,t} \cdot d_{l,t} = 0 \\
        &\forall l,\, t: && \omega^{d,\max}_{l,t} \cdot (d_{l,t} - D_{l,t}) = 0
    \end{align}

    \item \textbf{Factibilidad dual:}
    Las restricciones de desigualdad deben ser no negativas para encontrarse en el óptimo. Si alguna fuera negativa significaría que todavía se puede mejorar el valor de la función objetivo. Que todas deban ser no negativas pese a algunas ser cotas inferiores y otras superiores es porque las cotas inferiores ($0 \le x$) se pueden expresar como $-x \le 0$, lo cual resulta en un multplicador positivo.
    \begin{align}
        &\forall j,\, t: && \omega^{q,\min}_{j,t} \geq 0,\quad \omega^{q,\max}_{j,t} \geq 0 \\
        &\forall l,\, t: && \omega^{d,\min}_{l,t} \geq 0,\quad \omega^{d,\max}_{l,t} \geq 0
    \end{align}

    \item \textbf{Condiciones de estacionalidad:}
    Por último, para encontrarnos en el óptimo las derivadas parciales del Lagrangiano respecto a las variables del problema inferior deben ser 0. Estas variables son: $q_{i,t}, q_{j,t}, d_{l,t}$.
    \begin{align}
        &\forall j \neq i,\, \forall t: && -\alpha_{j,t} + \lambda_t - \omega^{q,\min}_{j,t} + \omega^{q,\max}_{j,t} = 0 \\
        &\forall t: && -\alpha_{i,t} + \lambda_t - \omega^{q,\min}_{i,t} + \omega^{q,\max}_{i,t} = 0 \\
        &\forall l,\, \forall t: && \beta_{l,t} - \lambda_t - \omega^{d,\min}_{l,t} + \omega^{d,\max}_{l,t} = 0
    \end{align}

\end{itemize}

\paragraph{}
\textbf{Linealización de las condiciones de complementariedad: big-M}

Para poder resolver el problema superior como un problema de optimización estándar, es necesario transformar las condiciones de complementariedad en restricciones lineales. Esto se logra mediante la técnica del big-M, introduciendo variables binarias y un parámetro $M$ suficientemente grande.

Para la condición de complementariedad:
\[
\omega^{q,\min}_{j,t} \cdot q_{j,t} = 0
\]
se sustituye por:
\begin{align}
    &\omega^{q,\min}_{j,t} \leq M \cdot z^{q,\min}_{j,t} 
    \label{bigm_1}\\
    &q_{j,t} \leq M \cdot (1 - z^{q,\min}_{j,t}) 
    \label{bigm_2}\\
    &z^{q,\min}_{j,t} \in \{0,1\}
\end{align}
donde $z^{q,\min}_{j,t}$ es una variable binaria y $M$ es el parámetro grande. Este esquema se aplica a todas las condiciones de complementariedad del modelo.

\paragraph{}
\textbf{Elección del parámetro $M$}

El valor de $M$ debe ser suficientemente grande para no restringir el espacio factible, pero no excesivamente grande para evitar problemas numéricos. 
Por ejemplo:
\begin{itemize}
    \item Usar como referencia las cotas naturales de las variables del modelo (por ejemplo, $Q_{j,t}$ y $D_{l,t}$) para \ref{bigm_2}. (Es curioso que esta haría que la cota no se rompa, pero como necesitamos la dual de esa restriccion para \ref{bigm_1} no se puede quitar la original)
    \item Para \ref{bigm_1} es más difícil elegir un M ajustado, pues en principio no sabemos que valores puede tomar la variable dual asocidada a la restricción. Si lo pensamos esto es cuanto puede afectar al Social Welfare que un productor no pueda producir más o menos. Es decir estas variables nos dicen el ahorro (gasto para las de min) que le daría a la sociedad si pudiera producir más. En el caso de min en principio no tiene mucho sentido, porque es lo que está diciendo es que si esta persona pudiera consumir en vez de vender. Es decir si estas en el limite inferior de un generador es porque el precio es demasiado alto, asi que querrias que compre en vez de vender. Me da la sensación que estas variables podrían ser acotadas por el precio máximo (de compra o venta) de la subasta.
\end{itemize}

\paragraph{}
\textbf{Linealización completa de las condiciones de complementariedad:}

\begin{itemize}
    \item \textbf{Generadores $j \neq i$:}
    \begin{align}
        &\omega^{q,\min}_{j,t} \leq M \cdot z^{q,\min}_{j,t} \\
        &q_{j,t} \leq M \cdot (1 - z^{q,\min}_{j,t}) \\
        &\omega^{q,\max}_{j,t} \leq M \cdot z^{q,\max}_{j,t} \\
        &q_{j,t} - Q_{j,t} \leq M \cdot (1 - z^{q,\max}_{j,t}) \\
        &z^{q,\min}_{j,t},\; z^{q,\max}_{j,t} \in \{0,1\}
    \end{align}
    \item \textbf{Generador estratégico $i$:}
    \begin{align}
        &\omega^{q,\min}_{i,t} \leq M \cdot z^{q,\min}_{i,t} \\
        &q_{i,t} \leq M \cdot (1 - z^{q,\min}_{i,t}) \\
        &\omega^{q,\max}_{i,t} \leq M \cdot z^{q,\max}_{i,t} \\
        &q_{i,t} - (1-x_{i,y(t)}) Q_{i,t} \leq M \cdot (1 - z^{q,\max}_{i,t}) \\
        &z^{q,\min}_{i,t},\; z^{q,\max}_{i,t} \in \{0,1\}
    \end{align}
    \item \textbf{Demanda $l$:}
    \begin{align}
        &\omega^{d,\min}_{l,t} \leq M \cdot z^{d,\min}_{l,t} \\
        &d_{l,t} \leq M \cdot (1 - z^{d,\min}_{l,t}) \\
        &\omega^{d,\max}_{l,t} \leq M \cdot z^{d,\max}_{l,t} \\
        &d_{l,t} - D_{l,t} \leq M \cdot (1 - z^{d,\max}_{l,t}) \\
        &z^{d,\min}_{l,t},\; z^{d,\max}_{l,t} \in \{0,1\}
    \end{align}
\end{itemize}

\section{Nivel superior: agente estratégico renovable}

El agente $i$, con un 20\% de la producción renovable total, decide anualmente la fracción $x_{i,y}$ a contratos PPAs. El resto, $(1-x_{i,y})\,Q_{i,t}$, se destina al mercado spot.

Problema del nivel superior:

\begin{align}
\max_{x}\;\;
& \sum_{t \in T} \lambda_t \Big( q_{i,t}\Big)
+ \sum_{t \in T} \gamma \Big(x_{i,y(t)} Q_{i,t} \Big)
\;-\; \sum_{t \in T} \alpha_{i,t}\, (q_{i,t} + x_{i,y(t)} Q_{i,t})
\label{UP_obj}\\[2mm]
\text{s.a.}\;\;
& 0 \le x_{i,y} \le 1 \quad \forall y \in Y
\label{UP_bounds}\\
& \text{KKT del nivel inferior en su versión lineal.}
\label{UP_MPEC}
\end{align}

% \subsection{Notas}

% \begin{itemize}
% \item \textbf{Estructura del modelo}
% \begin{itemize}
% \item Nivel superior: un agente estratégico $i$ con 20\% de la producción renovable del sistema decide anualmente qué fracción de su producción retira del spot vía PPA para maximizar su beneficio.
% \item Nivel inferior: \textit{market clearing} competitivo del mercado spot.
% \end{itemize}

% \item \textbf{Hipótesis clave}
% \begin{itemize}
% \item El agente $i$ solo posee generación renovable.
% \item Los PPAs no alteran la demanda spot. Esto hay que justificarlo, podemos decir que ante la electrificación de la industria aparecen oportunidades de PPAs que anteriormente no pasaban por mercado.
% \item La decisión del PPA se modela como un parámetro anual $x_{i,y}$.
% \end{itemize}

% \item \textbf{Resultados esperados}
% \begin{itemize}
% \item Retirar energía al PPA reduce la oferta del spot y puede elevar el precio $\lambda_t$.
% \item Ingresos = ventas en spot + ventas en PPA 
% \item Si el 100\% es PPA no se capturan los ingresos resultado de subir el precio.
% \end{itemize}

% \item \textbf{Caso A: PPA porcentaje constante anual}
% \begin{itemize}
% \item Decisión: $x_{i,y}$ fija; entregar $q^{PPA}_{t}=x_{i,y} Q_{i,t}$.
% \item Riesgo: el comprador asume variabilidad de volumen.
% \item Sentido: clientes flexibles con almacenamiento o procesos modulables que buscan certificación renovable.
% \item He escrito la formulación para este caso, puedes empezar por este o el siguiente como prefieras.
% \end{itemize}

% \item \textbf{Caso B: PPA perfil }
% \begin{itemize}
% \item Decisión: $\phi_y$ constante cada hora. Ahora sería una binaria que activa o no un perfil.
% \item La complicación añadida de este modelo es que hay que decidir los perfiles a ofertar, habría que buscar consumos industriales e inventarse algo parecido. Además Tendríamos que pensar si crear varios que sumándolos puedan llegar al 100\% de la capacidad de generación o solo uno y $\phi$ vaya de 0 a 1. 
% \item Al crear los perfiles habría que asegurarse de que siempre se pueden satisfacer, poner un coste de penalización cuando no se pueda (podría ser $\lambda$ asumiendo que lo compra en el mercado eléctrico).
% \item De cara a explicar los resultados el caso A me gusta más
% \item Riesgo: el vendedor asume el riesgo de producción.
% \item Sentido: clientes inflexibles que priorizan volumen estable y cobertura de precio.
% \end{itemize}



% \item \textbf{Perspectiva temporal}
% \begin{itemize}
% \item Los PPAs duran de 5 a 15 años o incluso más. Así que necesitamos un horizonte temporal suficientemente largo. Podemos empezar ejecutando de 2026 a 2036 y duración del PPA de un año. Si vamos bien en algún momento querremos activar PPAs de más años, para esto puedes usar las restricciones de tu formulación necesarias o usar binarias.
% \item En comparación con tu formulación un cambio importante es que los PPAs tienen duraciones de año natural, con esto se eliminan muchas variables y se resuelve más rápido. 
% \end{itemize}

% \item \textbf{Sobre tu formulación}
% \begin{itemize}
% \item Está bien, pero en modelos como este que son difíciles de resolver hay que simplificar. Por ahora, hemos simplificado cuándo empiezan los contratos (a 1 de enero) y su duración (año natural). Además usando variables anuales para x frente a variables horarias para q. Piensa cuántas variables y restricciones te quita esto por año para contarlo en la memoria. 
% \item Para la nomenclatura vamos a intentar utilizar nombres de una letra cuando sea posible, parte de la de este documento 
% \end{itemize}

% \item \textbf{Siguientes pasos y otras decisiones}
% \begin{itemize}
% \item Primero acaba el anexo B.
% \item Hay que terminar de definir el problema. El modelo de este documento está muy simplificado. Es suficiente para un TFG y nos va a permitir hacer cosas que con modelos más complejos no podríamos. Conviene que pienses en qué cambios podrían hacerse y por qué. Haz una lista y los discutimos. Algunos los integraremos y otros los contarás en la memoria. 
% \item Luego programa este modelo por KKTs o con algún algoritmo iterativo. (Habrá que hacer los dos pero para terminar de definir el problema empieza por el que quieras)
% \end{itemize}


% \end{itemize}

\end{document}
